\documentclass[11pt]{scrartcl}
\usepackage[a4paper,margin=2.6cm]{geometry}
\usepackage{../common/rlstudy}

\title{Extended reading --- Week 4: TD, SARSA, Q-learning, Expected SARSA}
\author{STUDY-RL}
\date{\today}
\begin{document}
\maketitle

\section*{Reading list}

The papers below are the original sources for the algorithms derived in
\texttt{notes.pdf}. For each: full citation, a one-paragraph summary, one annotated key
passage, and the curriculum thread it tightens.

\subsection*{1. Sutton 1988 --- ``Learning to Predict by the Methods of Temporal Differences''}

\textbf{Citation.} Richard~S.\ Sutton, \emph{Machine Learning} 3:9--44, 1988.

\paragraph{Summary.} The paper that defined TD methods and proved (informally) their
convergence under linear function approximation. Sutton frames TD as the natural
compromise between supervised-learning regression on Monte Carlo returns and the
dynamic-programming backup. He introduces the $\lambda$-return geometric mixture and
shows TD($\lambda$) is strictly more sample-efficient than MC on a 5-state random walk.

\paragraph{Why it matters for this week.} Every TD algorithm we will see for the next 20
weeks is a descendant of equation~(2) of this paper.

\begin{quote}\itshape
``Temporal-difference methods are a class of incremental learning procedures specialized
for prediction problems. Whereas conventional prediction-learning methods are driven by
the error between predicted and actual outcomes, TD methods are similarly driven by the
error or difference between temporally successive predictions; with them, learning occurs
whenever there is a change in prediction over time.''
\end{quote}
\textbf{Commentary.} Note Sutton's emphasis on \emph{change in prediction} as the learning
signal. That framing makes TD natural for online deployment: you do not need ``ground
truth,'' you only need a stream of predictions you trust to drift toward correctness.
Production systems often have exactly this structure --- noisy partial outcomes and a
stable online estimator.

\subsection*{2. Watkins \& Dayan 1992 --- ``Q-learning''}

\textbf{Citation.} Christopher~J.~C.~H.\ Watkins and Peter Dayan, \emph{Machine Learning}
8:279--292, 1992.

\paragraph{Summary.} Defines Q-learning and gives the first convergence proof for
off-policy control with a model-free agent. The argument constructs an action-value
process driven by a contraction in $\ell_\infty$ and verifies stochastic-approximation
conditions; the behaviour policy can be \emph{anything} that visits each $(s,a)$
infinitely often.

\paragraph{Why it matters for this week.} The clearest justification for why we can learn
the optimal policy by following a noisy exploratory policy. Foundation for everything in
Weeks 5--6 (DQN family).

\begin{quote}\itshape
``Q-learning works by directly approximating the action-value function $Q^\star$ ---
without ever requiring the agent to learn a model of the world or to evaluate alternative
policies.''
\end{quote}
\textbf{Commentary.} The structural separation between behaviour (what generates data) and
target (what we evaluate) is the central abstraction. The deadly triad of Week~5
(off-policy + bootstrapping + function approximation) is born from extending this
separation into the function-approximation regime where the convergence proof breaks.

\subsection*{3. van Hasselt 2010 --- ``Double Q-learning''}

\textbf{Citation.} Hado van Hasselt, \emph{NIPS} 2010.

\paragraph{Summary.} Identifies and analyses Q-learning's maximisation bias and proposes
maintaining two independently updated Q-estimators. Theorem~1 shows that
$\mathbb{E}[\max_a \hat Q(s,a)] \ge \max_a \mathbb{E}[\hat Q(s,a)]$ strictly when there is
noise across actions; the gap is upper bounded by the standard deviation of the noise.
Empirically demonstrates the fix on a roulette MDP.

\paragraph{Why it matters for this week.} Makes the maximisation bias precise. Double-DQN
in Week~6 is the function-approximation analogue.

\begin{quote}\itshape
``We have shown that the well-known Q-learning algorithm uses the $\max$ operator to
estimate the value of the maximum expected action value, leading to over-estimations of
the action values. Double Q-learning, which we introduce in this paper, uses two estimates
and decouples the selection from the evaluation, which results in an unbiased value
estimator.''
\end{quote}
\textbf{Commentary.} The decoupling principle (argmax from one estimator, value from
another) reappears in target networks (DQN), twin critics (TD3), and in ensemble
disagreement bonuses for exploration. Master it now.

\subsection*{4. van Seijen, van Hasselt, Whiteson, Wiering 2009 --- ``A Theoretical and Empirical Analysis of Expected Sarsa''}

\textbf{Citation.} Harm van Seijen, Hado van Hasselt, Shimon Whiteson, Marco Wiering,
\emph{IEEE ADPRL} 2009.

\paragraph{Summary.} Defines Expected SARSA, proves it has strictly lower variance than
SARSA (same expectation, integrate out next-action noise), and reports faster convergence
on cliff-walking-style tasks. Also discusses on- and off-policy variants: when the
expectation is taken under the behaviour policy you get on-policy ES; when taken under a
target policy you get a different off-policy estimator that subsumes Q-learning (target =
greedy) as a special case.

\paragraph{Why it matters for this week.} The cleanest example of a free-lunch variance
reduction in TD control; the algorithm also unifies SARSA and Q-learning under a single
objective.

\begin{quote}\itshape
``Expected Sarsa has the same convergence properties as Sarsa, but reduces the variance
in the updates and as a result performs better in practice. We argue that this variance
reduction is the main reason for the improved performance, and not a difference in bias.''
\end{quote}
\textbf{Commentary.} A nice example of Rao--Blackwellisation in RL: replace a sampled
estimator by its conditional expectation under a known distribution; never lose bias,
always reduce variance. Used implicitly in Expected-SARSA, in soft-Q targets, in GRPO's
group-mean baseline, and in many other algorithms we will see later.

\subsection*{5. Sutton \& Barto 2018 --- \emph{Reinforcement Learning: An Introduction}, Chapter 6}

\textbf{Citation.} Richard S.\ Sutton and Andrew G.\ Barto, MIT Press, 2nd edition, 2018.

\paragraph{Summary.} The textbook chapter that introduces TD prediction and control. The
cliff-walking example (Example 6.6) and the maximisation-bias example (Example 6.7) are
taken directly from here.

\paragraph{Why it matters for this week.} The standard reference; we re-derive its proofs
but use its problems verbatim.

\begin{quote}\itshape
``TD methods learn a guess from a guess --- they bootstrap.''
\end{quote}
\textbf{Commentary.} Three words that contain the whole of Weeks 4--10: ``learn a guess
from a guess'' is the entire game. The cost is the deadly triad of Week~5; the benefit is
online incremental learning that scales to billions of frames in Week~6.

\bigskip
\textit{Sourced from:} \texttt{research/deeprl\_core\_sota.md} \S1 (DQN family lineage),
and the canonical Sutton \& Barto textbook for the foundational material.

\end{document}
